\chapter*{Suggested Reading}
\addcontentsline{toc}{chapter}{Suggested Reading}
\markboth{Suggested Reading}{Suggested Reading}

This atlas unifies the dihedral geometry of composition mosaics across the elementary cyclic lattices $\mathbb{Z}_1$ through $\mathbb{Z}_8$, bridging integer partitions, circular set partitions, and dihedral equivalence classes. The curated literature below provides foundational theory, classical historical developments, and broader combinatorial and algebraic context.

\section*{Dihedral Enumeration and Orbit Counting}

\begin{description}
    \item[\textbf{Augustin-Louis Cauchy (1845)}]\mbox{}\\
    \emph{M\'emoire sur diverses propri\'et\'es remarquables des substitutions r\'eguli\`eres ou irr\'eguli\`eres, et des syst\`emes de substitutions conjugu\'ees}, \emph{Comptes Rendus de l'Acad\'emie des Sciences}, 21:972--987.\\
    The original formulation of the orbit-counting lemma for permutation groups.

    \item[\textbf{Ferdinand Georg Frobenius (1887)}]\mbox{}\\
    \emph{\"Uber die Congruenz nach einem aus zwei endlichen Gruppen gebildeten Doppelmodul}, \emph{Journal f\"ur die reine und angewandte Mathematik (Crelle's Journal)}, 101:273--299.\\
    The generalisation of the orbit-counting theorem to abstract finite group actions.

    \item[\textbf{William Burnside (1911)}]\mbox{}\\
    \emph{Theory of Groups of Finite Order}, 2nd edition, Cambridge University Press, Cambridge.\\
    The classical treatise popularising the orbit-counting lemma, forming the arithmetic core of dihedral cycle index evaluations on cyclic lattices.

    \item[\textbf{George P\'olya (1937)}]\mbox{}\\
    \emph{Kombinatorische Anzahlbestimmungen f\"ur Gruppen, Graphen und chemische Verbindungen}, \emph{Acta Mathematica}, 68(1):145--254.\\
    The foundational memoir formalising the cycle index polynomial $Z(G; z_1, \dots, z_n)$ and its functional substitution principles for weighted configuration counting under group equivalence.

    \item[\textbf{Harald Fripertinger (1997)}]\mbox{}\\
    \emph{Enumeration of isometry classes of mosaics}, \emph{Discrete Mathematics}, 176(1--3):87--98.\\
    The pioneering combinatorial census of set partitions on cyclic and dihedral lattices, establishing the orbit inventory of geometric mosaics.

    \item[\textbf{Adalbert Kerber (1999)}]\mbox{}\\
    \emph{Applied Finite Group Actions}, 2nd edition, Springer-Verlag, Berlin, Heidelberg.\\
    An exhaustive modern reference on the algebraic apparatus of permutation representations, double-coset decompositions, and constructive orbit generation.
\end{description}

\section*{Set Partitions, Combinatorial Species, and Fa\`a di Bruno Theory}

\begin{description}
    \item[\textbf{Eric Temple Bell (1934)}]\mbox{}\\
    \emph{Exponential polynomials}, \emph{Annals of Mathematics}, 35(2):258--277.\\
    Foundational formulation of exponential polynomials and Bell numbers enumerating set partitions, establishing the combinatorial bedrock of total mosaic volume across $\mathbb{Z}_1 \dots \mathbb{Z}_8$.

    \item[\textbf{Gian-Carlo Rota (1964)}]\mbox{}\\
    \emph{The Number of Partitions of a Set}, \emph{The American Mathematical Monthly}, 71(5):498--504.\\
    The classic modern foundation of partition lattice combinatorics, generating functions, and umbral operator techniques for set partition enumerations.

    \item[\textbf{Andr\'e Joyal (1981)}]\mbox{}\\
    \emph{Une th\'eorie combinatoire des s\'eries formelles}, \emph{Advances in Mathematics}, 42(1):1--82.\\
    The categorical foundation of combinatorial species, providing the functorial machinery underlying Singleton Shattering Duality and cycle index series.

    \item[\textbf{Francesco Fa\`a di Bruno (1855)}]\mbox{}\\
    \emph{Sullo sviluppo delle funzioni}, \emph{Annali di Scienze Matematiche e Fisiche}, 6:479--480.\\
    The classical multinomial expansion governing the derivative of composite functions, structuring the shape-stratified multinomial decomposition of circular compositions.

    \item[\textbf{Philippe Flajolet and Robert Sedgewick (2009)}]\mbox{}\\
    \emph{Analytic Combinatorics}, Cambridge University Press, Cambridge.\\
    The definitive modern treatise on combinatorial constructions, exponential generating functions, and multiset formations under symmetry.
\end{description}

\section*{Difference Systems and Combinatorial Block Designs}

\begin{description}
    \item[\textbf{Thoralf Skolem (1957)}]\mbox{}\\
    \emph{On certain distributions of integers in pairs with given differences}, \emph{Mathematica Scandinavica}, 5:57--68.\\
    The seminal paper introducing difference systems of integer pairs, providing the structural prototype for the unique cyclic Skolem mosaic partition on $\mathbb{Z}_8$.

    \item[\textbf{Duncan M. Y. Sommerville (1909)}]\mbox{}\\
    \emph{On Certain Periodic Properties of Cyclic Compositions of Numbers}, \emph{Proceedings of the London Mathematical Society}, Second Series, 7:263--313.\\
    Early 20th-century foundation for periodic cyclic integer compositions and circular stride geometries.
\end{description}

\section*{The Cyclic Lattice Cartography Trilogy}

\begin{description}
    \item[\textbf{Gemma Knee and Antigra Vitti (2026)}]\mbox{}\\
    \emph{The Polygonal Burnsidean: An Atlas of Inscribed Polygons from $\mathbb{Z}_1$ to $\mathbb{Z}_{16}$}, Newcastle upon Tyne.\\
    The first volume in the trilogy: complete algebraic classification and visual cartography of all 5,099 dihedral orbits of inscribed single polygons across $\mathbb{Z}_1 \dots \mathbb{Z}_{16}$.

    \item[\textbf{Gemma Knee and Antigra Vitti (2026)}]\mbox{}\\
    \emph{Hexadecality: The Atlas of Hexadecagonal Compositions}, Newcastle upon Tyne.\\
    The second volume: exhaustive affine macro-orbit classification and visual plates of the 21,030 macro-orbits (83,488 dihedral comosaics) of uniform tetramers ($z_4^4$) on $\mathbb{Z}_{16}$, alongside the square universes of $\mathbb{Z}_4$ and $\mathbb{Z}_9$.
\end{description}
